\appendix
\section{Corpus composition}
\label{app:corpus}

Table~\ref{tab:corpus} summarises the corpus structure and lists one sample
sentence per stratum. The four domains are constructed to differ in register,
vocabulary, and numerical density. The broadcast-news domain uses formal
written register with institutional vocabulary, the conversation domain uses
colloquial everyday phrasing, the technical domain uses scientific and
engineering terminology, and the numerals-and-entities domain concentrates
dates, times, quantities, prices, addresses, codes, and named entities, with
all numbers written out in words as described in
Section~\ref{sec:methods}. Length strata are defined per language as short,
medium, and long as given in Section~\ref{sec:methods}, and every
domain-by-length cell contains twenty sentences per language, giving 240
sentences per language and 480 in total. The severity sweep subset of 144
sentences is sampled evenly across all cells, six sentences per cell per
language.

\begin{table}[h]
\centering
\caption{Corpus structure with one sample sentence per domain.}
\label{tab:corpus}
\small
\begin{tabular}{lp{8.8cm}}
\toprule
Domain & Sample sentence \\
\midrule
News (zh) & 国家统计局今天发布了上月经济运行数据。 \\
News (en) & The central bank raised interest rates by twenty five basis points to curb inflation. \\
Conversation (zh) & 我打算下班之后去健身房锻炼一个小时。 \\
Conversation (en) & I was thinking we could drive to the beach this weekend if the weather holds. \\
Technical (zh) & 混响是声波在房间内多次反射形成的，会显著降低远场语音识别的准确率。 \\
Technical (en) & Reverberation arises when sound waves reflect repeatedly off room surfaces. \\
Numerals (zh) & 从北京到上海的高铁大约需要四个半小时。 \\
Numerals (en) & The high speed train from Beijing to Shanghai takes about four and a half hours. \\
\bottomrule
\end{tabular}
\end{table}

\section{Qualitative error analysis}
\label{app:qualitative}

Table~\ref{tab:examples} collects representative model responses that
illustrate the failure patterns behind the aggregate numbers in
Section~\ref{sec:results}.

The refusal behaviour of Section~\ref{sec:results:main} is shown in the
first row: under 0~dB babble the Qwen-Audio recogniser returns an empty
transcript for a substantial English sentence, which the benchmark scores as
a complete failure. The second row shows the opposite regime: the same
system, on a long English news sentence at the same severity, produces a
verbatim transcript, which underlines how bimodal the sentence-level
distributions of Fig.~\ref{fig:dist} are. Row three shows the dominant error
mode of the specialised streaming recogniser under severe noise, where most
of the sentence is lost and only a short fragment survives. Row four shows a
typical near-miss of the LLM recogniser in Chinese, where a single
content word, density, is replaced by another word of similar phonetic
shape, speed, without any structural damage. Row five shows the collapse of
the telephone-band specialist under reverberation, which occasionally
produces short disconnected fragments rather than transcripts. Row six
shows an error that the recognition stage never sees but the translation
consumer experiences: under 0~dB noise the omni end-to-end path substitutes
a plausible but wrong entity, rendering the sports centre of the source as a
grand theatre, a hallucination that a string-matching metric penalises only
lightly but that is catastrophic for factual fidelity. Row seven shows the
narrowband immunity of the LLM recogniser, which returns a verbatim
transcript of telephone-band audio. Together the examples support the
reading of Section~\ref{sec:discussion}: the systems differ less in whether
they fail than in how, and the failure modes, from verbatim robustness to
silent entity substitution to outright refusal, have distinct downstream
consequences.

\begin{table}[h]
\centering
\caption{Representative system responses under degradation.}
\label{tab:examples}
\small
\begin{tabular}{llp{6.1cm}}
\toprule
Condition & System & Reference and response \\
\midrule
0 dB, en & Qwen-Audio-3.1 & Reference: Police arrested two suspects on Tuesday. Response: empty transcript. \\
0 dB, en & Qwen-Audio-3.1 & Reference: The transport minister told reporters on Friday that the government would invest heavily in electric vehicle charging infrastructure. Response: verbatim transcript. \\
0 dB, en & Paraformer-RT-v2 & Reference: The battery lasts twenty hours. Response: About. \\
0 dB, zh & Qwen-Audio-3.1 & Reference: 这款处理器的晶体管密度提升了近一倍。 Response: 这款处理器的晶体管速度提升了近一倍。 \\
Reverb, zh & Paraformer-RT-8K & Reference: 记者从发布会上了解到，本次合作项目将覆盖全国二十多个主要城市。 Response: 七月同。非常了解到。 \\
0 dB, E2E & Omni-Flash & Source: 新建成的体育中心可容纳六万名观众。 Translation: The newly built Grand Theatre can accommodate 60,000 audience members. \\
Narrow, en & Qwen-Audio-3.1 & Reference: This coffee shop is cozy. Response: verbatim transcript. \\
$-10$ dB, en & Qwen-Audio-3.1 & Reference: The battery lasts twenty hours. Response: fluent unrelated Chinese text of several sentences (hallucination). \\
\bottomrule
\end{tabular}
\end{table}

\section{Cross-system text normalisation}
\label{app:normalisation}

Comparing transcripts across commercial recognisers requires a normalisation
protocol that neutralises provider-specific inverse text normalisation. The
protocol applies the following rules identically to references and
hypotheses, in order. Unicode compatibility folding and case folding are
applied first, and punctuation is removed. Numeric expressions are then
canonicalised to digit form on both sides. For Chinese, the rules handle
cardinal numbers with the standard multipliers, digit-by-digit readings such
as room and code numbers where each digit is read separately, decimal
fractions written with the point particle, time-of-day expressions in which
the hour and minute parts are converted separately, percentage markers,
which are deleted on both sides, and scale multipliers, where a digit
followed by wan, yi, or their combinations is expanded to the full
multiplied-out form, so that twenty-six million written as 2600万 in one
system's output and as 26000000 in another's both canonicalise identically.
For English, the rules handle cardinal numbers including compound
hyphenated forms and the and connector, time and digit-by-digit readings
such as twelve zero eight and nine forty, which canonicalise to 1208 and
940, decimal fractions expressed with point, scale multipliers from thousand
to trillion, and ordinal expressions of both word and digit form, which are
masked to a common placeholder since their spoken-to-written realisations
differ across systems in ways that carry no acoustic information. The rule
set was developed and frozen by inspecting all high-error sentences of the
clean condition until no residual mismatches attributable to normalisation
remained, a procedure that could in principle favour whichever system's
conventions dominate the clean outputs, so a sensitivity analysis was run:
re-scoring the main grid with case and punctuation folding only, dropping
numeral canonicalisation entirely, raises clean-condition error rates from
0.5, 2.3, and 0.3 percent to 6.1, 6.2, and 0.8 percent for Qwen-Audio,
Paraformer, and Qwen3-ASR respectively, while the 0~dB scores move by less
than 1.5 points and the system ordering is unchanged in every condition.
The normalisation therefore compresses but does not create the reported
system differences, and the residual protocol noise on clean speech is
bounded above by the strongest system's per-language clean rate of 0.14
percent in Mandarin (0.29 percent pooled over both languages), with the
caveat that the two systems most affected by numeral conventions differ
most under the raw protocol. The full normalisation code is released with
the benchmark.

